\documentclass[11pt]{article}
\usepackage[a4paper,margin=1in]{geometry}
\usepackage{amsmath,amssymb,amsthm,booktabs,enumitem,mathtools}
\usepackage[T1]{fontenc}
\usepackage{lmodern}
\usepackage{hyperref}

\hypersetup{colorlinks=true,linkcolor=blue,citecolor=blue,urlcolor=blue}
\newtheorem{theorem}{Theorem}[section]
\newtheorem{proposition}[theorem]{Proposition}
\newtheorem{lemma}[theorem]{Lemma}
\newtheorem{corollary}[theorem]{Corollary}
\theoremstyle{definition}
\newtheorem{definition}[theorem]{Definition}
\theoremstyle{remark}
\newtheorem{remark}[theorem]{Remark}
\newcommand{\Mob}{\mu}
\newcommand{\one}{\mathbf 1}
\newcommand{\E}{\mathbb E}
\newcommand{\Pp}{\mathbb P}

\title{Mixed-exponent M\"obius run hierarchies and
the two-envelope capacity boundary}
\author{RH research program}
\date{August 2026}

\begin{document}
\maketitle

\begin{abstract}
The exact RH-371 distance-two capacity uses sixteen overlapping signed-run
counts, two signs at each length $1\le k\le8$.  We expand every count on its
native odd-start endpoint into mixed M\"obius exponents.  If $H_{k,r}$ is the
sum of all terms with $r$ first powers and $k-r$ square masks, and $A_k,B_k$
collect the layers $r\ge2$ of even and odd parity, then
\[
 2^kC_{\sigma,k}=H_{k,0}+A_k+
 \sigma(H_{k,1}+B_k).
\]
We prove unconditionally that
$H_{k,0}/N\to\Delta_k=\frac12\prod_{p\ \mathrm{odd}}(1-k/p^2)$ and
$H_{k,1}=o(N)$.  The latter argument uses bounded periodic prime-square
masks at fixed cutoff, fixed-progression Davenport cancellation, and only
then a large-prime-square union bound.  Consequently all sixteen densities
exist simultaneously exactly when thirteen aggregate limits exist:
$A_k/N$ for $2\le k\le8$ and $B_k/N$ for $3\le k\le8$.  The associated
formal $466$-coordinate block map has rank $13$ and kernel dimension $453$.
This is not an arithmetic minimality theorem.  For the adaptive capacity we
obtain
\[
 K_N/N=2r_0+2\{U_N+|V_N|\}/N+o(1),
\]
so convergence is equivalent to one two-envelope limit.  Full
mixed-exponent cancellation is sufficient, but not necessary, for the
displayed conditional Euler-product constant.  A self-contained stationary
ternary chain shows that raw, square-only, and one-sign masked moments do
not algebraically determine the first directional two-sign masked moment.
The chain is synthetic and not M\"obius.  No capacity limit, operator,
trace formula, zero identification, Hilbert--Polya construction, or proof
of RH is claimed.
\end{abstract}

\section{Native endpoints and mixed-exponent coordinates}

The frozen RH-371 reduction writes the distance-two adaptive capacity in
terms of signed intervals at lengths at most eight \cite{RH371}.  For
$I_k=\{0,\ldots,k-1\}$ set
\begin{equation}
 \mathcal D_k(N)=
 \{n\in\mathbb Z:1\le n\le N-2(k-1),\ n\text{ odd}\}.
 \label{eq:domain}
\end{equation}
Every object at level $k$ below uses exactly this endpoint.  Empty domains
give zero.  For $S\subseteq I_k$, define
\begin{equation}
 T_{k,S}(N)=\sum_{n\in\mathcal D_k(N)}
 \prod_{j\in S}\Mob(n+2j)
 \prod_{j\in I_k\setminus S}\Mob(n+2j)^2,
 \qquad
 H_{k,r}(N)=\sum_{\substack{S\subseteq I_k\\|S|=r}}T_{k,S}(N).
 \label{eq:mixed-moments}
\end{equation}
The signed interval count is
\begin{equation}
 C_{\sigma,k}(N)=\#\{n\in\mathcal D_k(N):
       \Mob(n)=\cdots=\Mob(n+2(k-1))=\sigma\},
 \quad \sigma\in\{-1,+1\}.
 \label{eq:run-count}
\end{equation}
These are overlapping same-sign windows.  They are not counts of maximal
runs or runs having exact length $k$.

Separate the higher mixed layers by parity:
\begin{equation}
 A_k=\sum_{\substack{2\le r\le k\\r\ \mathrm{even}}}H_{k,r},
 \qquad
 B_k=\sum_{\substack{3\le r\le k\\r\ \mathrm{odd}}}H_{k,r}.
 \label{eq:AB}
\end{equation}
Thus $A_1=B_1=B_2=0$.  The first nonzero aggregate is
$A_2=T_{2,\{0,1\}}$, the ordinary shift-two correlation on the odd-start
endpoint isolated in RH-376 \cite{RH376}.

\begin{proposition}[All-prefix mixed identity]
For every $N\ge1$, $1\le k\le8$, and $\sigma\in\{-1,+1\}$,
\begin{equation}
 \boxed{\displaystyle
 2^kC_{\sigma,k}=H_{k,0}+A_k+
                  \sigma\bigl(H_{k,1}+B_k\bigr).}
 \label{eq:master}
\end{equation}
\end{proposition}

\begin{proof}
For $x\in\{-1,0,1\}$,
$2\one_{\{x=\sigma\}}=x^2+\sigma x$.  Hence, pointwise in
$n\in\mathcal D_k(N)$,
\[
 2^k\prod_{j\in I_k}\one_{\{\Mob(n+2j)=\sigma\}}
 =\prod_{j\in I_k}\bigl(\Mob(n+2j)^2+\sigma\Mob(n+2j)\bigr).
\]
Expansion over $S\subseteq I_k$, followed by $\sigma^{|S|}=1$ for even
$|S|$ and $\sigma$ for odd $|S|$, gives \eqref{eq:master} after summation.
No endpoint extension or probabilistic assumption occurs.
\end{proof}

\section{The two deterministic layers}

Put
\begin{equation}
 e_k=\prod_{p\ \mathrm{odd}}\left(1-\frac{k}{p^2}\right),
 \qquad \Delta_k=\frac{e_k}{2},
 \qquad 1\le k\le8.
 \label{eq:delta}
\end{equation}
The product is only over odd primes.  The factor $1/2$ outside the product
is exactly the odd-start condition.  No $p=2$ local factor is inserted into
$e_k$.

\begin{theorem}[Zero-sign and one-sign layers]
For each fixed $1\le k\le8$,
\begin{equation}
 \frac{H_{k,0}(N)}{N}\longrightarrow\Delta_k,
 \qquad H_{k,1}(N)=o(N).
 \label{eq:H01}
\end{equation}
\end{theorem}

\begin{proof}
For $H_{k,0}$, the summand is the indicator that all $k$ odd integers
$n,n+2,\ldots,n+2(k-1)$ are squarefree.  Fix a prime cutoff $R$.  For an
odd prime $p$, the $k$ shifts occupy distinct classes modulo $p^2$: if
$2(j-\ell)\equiv0\pmod{p^2}$, then $p^2\mid j-\ell$, which is impossible
for $0<|j-\ell|\le7<p^2$.  The Chinese remainder theorem therefore gives
the finite-sieve density
\[
 \frac12\prod_{\substack{p\le R\\p\ \mathrm{odd}}}
 \left(1-\frac{k}{p^2}\right).
\]
An accepted start omitted by the true squarefree condition has
$p^2\mid n+2j$ for some $j\in I_k$ and odd prime $p>R$.  The union bound is
\[
 O_k\!\left(N\sum_{p>R}\frac1{p^2}+\sqrt N\right)
 =O_k(N/R+\sqrt N).
\]
First send $N\to\infty$ at fixed $R$ and then $R\to\infty$.  This proves
the first limit, in the classical squarefree-pattern framework
\cite{Mirsky1948}.

For the second assertion, fix $j\in I_k$ and write the corresponding term
of $H_{k,1}$ as
\begin{equation}
 \sum_{n\in\mathcal D_k(N)}\Mob(n+2j)
       \prod_{\ell\ne j}\Mob(n+2\ell)^2.
 \label{eq:one-sign-term}
\end{equation}
For a fixed cutoff $R$, let
\[
 g_R(m)=\one_{\{p^2\nmid m\text{ for every prime }p\le R\}}.
\]
Replace each square mask in \eqref{eq:one-sign-term} by $g_R$.  The product
of these masks and the odd-start indicator is a bounded periodic function
with a fixed period depending only on $R$ and $k$.  Expanding it into its
finitely many residue classes turns the truncated sum, after translating
$m=n+2j$, into finitely many fixed-progressions sums of $\Mob(m)$.  Each is
$o(N)$ by Davenport's fixed-frequency theorem and finite Fourier inversion
\cite{Davenport1937}.

The true and truncated products can differ only if, for some
$\ell\ne j$, a prime $p>R$ satisfies $p^2\mid n+2\ell$.  Because the masks
are bounded by one, another union bound gives
\[
 O_k\!\left(N\sum_{p>R}\frac1{p^2}+\sqrt N\right)
 =O_k(N/R+\sqrt N).
\]
There are only $k$ possible sign positions.  Divide by $N$, take
$N\to\infty$ with $R$ fixed, and then take $R\to\infty$ to obtain
$H_{k,1}=o(N)$.
\end{proof}

\begin{remark}[Quantifier order]
The one-sign proof uses bounded periodic $0/1$ masks.  It does not expand
several unrestricted divisor sums and does not apply Davenport to a modulus
growing with $N$.  The order is fixed $R$, then $N\to\infty$, then
$R\to\infty$.
\end{remark}

For $k=2$, $e_2/2=\prod_p(1-2/p^2)$, because the missing $p=2$ factor of
the full product is $1/2$.  Thus the $k=2$ specialization agrees exactly
with the RH-376 squarefree-pair normalization.

\section{Thirteen aggregate density channels}

Dividing \eqref{eq:master} by $N$ and using \eqref{eq:H01} gives
\begin{equation}
 \frac{2^kC_{\sigma,k}(N)}N
 =\Delta_k+\frac{A_k(N)}N+\sigma\frac{B_k(N)}N+o(1).
 \label{eq:density-reduction}
\end{equation}

\begin{theorem}[Simultaneous density boundary]
All sixteen limits
\[
 \lim_{N\to\infty}C_{\sigma,k}(N)/N,
 \qquad \sigma\in\{-1,+1\},\quad1\le k\le8,
\]
exist if and only if all thirteen limits
\begin{equation}
 \lim_{N\to\infty}A_k(N)/N\quad(2\le k\le8),
 \qquad
 \lim_{N\to\infty}B_k(N)/N\quad(3\le k\le8)
 \label{eq:thirteen}
\end{equation}
exist.
\end{theorem}

\begin{proof}
If the limits in \eqref{eq:thirteen} exist, equation
\eqref{eq:density-reduction} gives all sixteen signed limits.  Conversely,
for each $k\ge3$, adding the $\sigma=+1$ and $\sigma=-1$ equations recovers
$A_k/N$ up to a convergent deterministic term and $o(1)$.  Subtracting them
recovers $B_k/N+o(1)$.  At $k=2$, $B_2=0$, so either signed equation
recovers $A_2/N$.  At $k=1$ there is no higher aggregate.  This yields
exactly seven $A$ limits and six $B$ limits.
\end{proof}

For a fixed $k\ge3$, convergence for only one sign controls the single
combination $(A_k+\sigma B_k)/N$.  It need not give separate convergence of
$A_k/N$ and $B_k/N$.

There are
\begin{equation}
 \sum_{k=2}^8(2^k-1-k)=466
 \label{eq:466}
\end{equation}
formal coordinates $T_{k,S}$ with $|S|\ge2$.  Consider the purely formal
linear map that sums even-cardinality coordinates into $A_k$ and
odd-cardinality coordinates into $B_k$.  Its thirteen row supports are
pairwise disjoint and nonempty.  Their sizes are
\[
 (1,3,7,15,31,63,127)\quad\text{and}\quad
 (1,4,11,26,57,120).
\]
Hence the map has rank $13$ and kernel dimension $466-13=453$.  This count
treats the coordinates as formal independent variables.  It does not prove
that thirteen is the minimal number of actual M\"obius correlation limits,
nor does it identify arithmetic relations among those limits.

\section{Exact two-envelope capacity reduction}

Let $E_\sigma(N)$ be the isolated even-path count from RH-371 and set
\begin{equation}
 s_k=(-1)^{k+1},\qquad
 R_\sigma(N)=E_\sigma(N)+\sum_{k=1}^8s_kC_{\sigma,k}(N).
 \label{eq:R}
\end{equation}
The locked path formula identifies $R_\sigma$ with a nonnegative
maximum-weight independent-set value and gives
\begin{equation}
 K_N=\max\{|{-M_N+2R_+(N)}|,|{-M_N-2R_-(N)}|\},
 \qquad M_N=\sum_{n\le N}\Mob(n).
 \label{eq:K-exact}
\end{equation}

Define four exact finite channels
\begin{align}
 P_N&=\frac{E_+(N)+E_-(N)}2+
       \sum_{k=1}^8\frac{s_kH_{k,0}(N)}{2^k},
 &Q_N&=\frac{E_+(N)-E_-(N)}2+
       \sum_{k=1}^8\frac{s_kH_{k,1}(N)}{2^k},
 \label{eq:PQ}\\
 U_N&=\sum_{k=2}^8\frac{s_kA_k(N)}{2^k},
 &V_N&=\sum_{k=3}^8\frac{s_kB_k(N)}{2^k}.
 \label{eq:UV}
\end{align}

\begin{proposition}[Finite channel and residual identities]
For every $N$ and $\sigma\in\{-1,+1\}$,
\begin{align}
 R_\sigma&=P_N+U_N+\sigma(Q_N+V_N),
 \label{eq:R-decomposition}\\
 \max(R_+,R_-)&=P_N+U_N+|Q_N+V_N|,
 \label{eq:maxR}\\
 \left|K_N-2\bigl(P_N+U_N+|V_N|\bigr)\right|
 &\le |M_N|+2|Q_N|.
 \label{eq:residual-bound}
\end{align}
\end{proposition}

\begin{proof}
Insert \eqref{eq:master} into \eqref{eq:R} and group the four types of
terms to obtain \eqref{eq:R-decomposition}.  Taking the larger sign gives
\eqref{eq:maxR}.  Since $R_\pm\ge0$, equation \eqref{eq:K-exact} and the
reverse triangle inequality give
\[
 |K_N-2\max(R_+,R_-)|\le|M_N|.
\]
Finally,
$\bigl||Q_N+V_N|-|V_N|\bigr|\le|Q_N|$, which proves
\eqref{eq:residual-bound}.
\end{proof}

The fixed-progression squarefree count on $n\equiv2\pmod4$ and Davenport
cancellation give
\[
 \frac{E_++E_-}{2N}\longrightarrow\frac1{\pi^2},
 \qquad E_+-E_-=o(N),\qquad M_N=o(N).
\]
Indeed, writing $n=2m$ identifies $E_++E_-$ with the number of odd
squarefree $m\le N/2$.  Their density among all integers is $4/\pi^2$, so
$E_++E_-=2N/\pi^2+o(N)$.  Moreover,
$E_+-E_-=\sum_{n\le N,\ n\equiv2\ (4)}\Mob(n)$, while $M_N$ is the
corresponding modulus-one sum.  Both are $o(N)$ by the same fixed-progression
Davenport input.
Together with \eqref{eq:H01}, these imply
\begin{equation}
 P_N=r_0N+o(N),\qquad Q_N=o(N),\qquad
 r_0=\frac1{\pi^2}+\sum_{k=1}^8\frac{s_k\Delta_k}{2^k}.
 \label{eq:r0}
\end{equation}

\begin{theorem}[Two-envelope capacity boundary]
Unconditionally,
\begin{align}
 R_\sigma(N)&=r_0N+U_N+\sigma V_N+o(N),
 \label{eq:R-asymptotic}\\
 \frac{K_N}{N}&=2r_0+
       2\frac{U_N+|V_N|}{N}+o(1).
 \label{eq:K-asymptotic}
\end{align}
Consequently $K_N/N$ converges if and only if
$(U_N+|V_N|)/N$ converges.
\end{theorem}

\begin{proof}
Equations \eqref{eq:R-asymptotic} and \eqref{eq:K-asymptotic} follow from
\eqref{eq:R-decomposition}, \eqref{eq:residual-bound}, and
\eqref{eq:r0}.  Since $r_0$ is fixed and the remaining error tends to zero,
the asserted equivalence follows in both directions.
\end{proof}

If every $T_{k,S}(N)=o(N)$ with $|S|\ge2$, then $U_N,V_N=o(N)$ and the
conditional capacity limit would be
\begin{equation}
 \boxed{\displaystyle
 \frac{2}{\pi^2}+\sum_{k=1}^8\frac{(-1)^{k+1}e_k}{2^k}.}
 \label{eq:conditional-constant}
\end{equation}
This full mixed-exponent cancellation hypothesis is sufficient only.  The
actual criterion is convergence of one envelope, so cancellation or
compensation among aggregates could suffice without every coordinate being
$o(N)$.  Neither condition is proved here.

\section{A stationary ternary algebraic boundary}

The following witness tests what the low layers can imply by stationarity
and ternary algebra alone.  It is not an arithmetic model for M\"obius.

\begin{proposition}[Directional second-order witness]
Fix $0<|\varepsilon|<3/2$.  There is a stationary process
$(X_t)_{t\in\mathbb Z}$ on $\{-1,0,1\}$ with uniform stationary pair law and
second-order transition
\begin{equation}
 \Pp(X_{t+1}=c\mid X_{t-1}=a,X_t=b)
 =\frac13\left[1+\varepsilon ab\left(c^2-\frac23\right)\right].
 \label{eq:transition}
\end{equation}
It has the following properties.
\begin{enumerate}[label=(\roman*),leftmargin=*]
\item every nonempty raw moment at distinct times is zero.
\item every square-only moment at $r$ distinct times equals $(2/3)^r$.
\item every moment with exactly one first-power factor and otherwise only
      square masks is zero.
\item for $A=X_0,B=X_1,C=X_2$,
\[
 \E[ABC^2]=\frac{8\varepsilon}{81},\qquad
 \E[AB^2C]=\E[A^2BC]=0.
\]
\item
$\Pp(A=B=C=1)=\Pp(A=B=C=-1)=
  27^{-1}(1+\varepsilon/3)$.
\end{enumerate}
\end{proposition}

\begin{proof}
For $ab\in\{-1,0,1\}$ and $c^2-2/3\in\{-2/3,1/3\}$, the bound on
$\varepsilon$ makes every probability in \eqref{eq:transition} positive.
Summing over $c$ gives one.  If $(A,B)$ is uniform, then for each $(b,c)$
\[
 \sum_a\frac19\,\Pp(c\mid a,b)=\frac19,
\]
because $\sum_aa=0$.  Hence the uniform pair law is stationary.

Directly from \eqref{eq:transition},
\begin{equation}
 \E[X_{t+1}\mid X_{t-1},X_t]=0,
 \qquad
 \E[X_{t+1}^2\mid X_{t-1},X_t]
 =\frac23+\frac{2\varepsilon}{9}X_{t-1}X_t.
 \label{eq:conditional-moments}
\end{equation}
For a raw product at arbitrary distinct times, condition on the history
immediately before the latest factor.  The first identity in
\eqref{eq:conditional-moments} gives zero.  For a square-only product,
condition at the latest squared factor.  The correction term contains
$X_{t-1}X_t$ times earlier squares.  Since $x^{2q+1}=x$ on the ternary
alphabet, conditioning on the latest of these two first-power factors again
annihilates the correction.  Induction leaves a factor $2/3$ for each
selected time.

The uniform pair law and transition are invariant under the global sign
flip $X\mapsto-X$.  A monomial with exactly one first power changes sign,
proving (iii).  The second identity in \eqref{eq:conditional-moments} gives
\[
 \E[ABC^2]=\frac{2\varepsilon}{9}\E[A^2B^2]
 =\frac{2\varepsilon}{9}\left(\frac23\right)^2
 =\frac{8\varepsilon}{81}.
\]
The other two directional moments vanish by conditioning on $C$.  Finally,
the probability of a fixed triple is the uniform pair mass $1/9$ times
\eqref{eq:transition}.  Substituting $(1,1,1)$ and $(-1,-1,-1)$ proves (v).
\end{proof}

At $\varepsilon=1$, the witness gives $\E[ABC^2]=8/81$ and both same-sign
triple probabilities $4/81$.  On the odd lattice its one-site nonzero
probability is $2/3$, whereas the M\"obius odd-lattice density is
$e_1=8/\pi^2$.  Equivalently, after embedding the chain at odd starts and
normalizing by the original $N$, the synthetic contribution is $1/3$,
versus $\Delta_1=4/\pi^2$.  Thus even its square layer does not match the
arithmetic density.  The proposition proves only that, in the
general stationary ternary class, the directional two-sign masked layer is
not forced by raw, square-only, and one-sign masked data.  It is neither a
M\"obius counterexample nor evidence that any M\"obius limit fails.

\section{Executable protocol and frozen rows}

The standard-library artifact has four independent components.  First, it
checks \eqref{eq:master} on all $19680$ ternary-window/sign cases for
$1\le k\le8$ and performs exact rational row reduction on the formal
$466\to13$ block map.  Second, an integer linear sieve streams every native
odd-start window through $N=2^{18}$.  It checks $1048548$ window updates,
$4194304$ cumulative signed identities, and all $262144$ prefix path-DP,
$R_\sigma$, $K_N$, and residual relations.  Third, exact rational arithmetic
checks the stationary chain: $27$ transitions, $9$ stationarity cells,
$502$ raw cases, $502$ square-only cases, and $1793$ one-sign masked cases
for block lengths at most eight.  The analytic proof above, not those finite
enumerations, establishes the arbitrary-distinct-time statements.  Fourth,
a finite Euler diagnostic uses all primes through $2^{20}$.

For compact display, the exact finite channels in the next table are
multiplied by $256$.  The final two columns are respectively the left side
and right side of the scaled version of \eqref{eq:residual-bound}.
\begin{center}
\small
\begin{tabular}{@{}rrrrrrr@{}}
\toprule
$N$ & $256P_N$ & $256Q_N$ & $256U_N$ & $256V_N$ & residual & bound\\
\midrule
$1024$ & $64711$ & $1108$ & $1465$ & $44$ & $3240$ & $3240$\\
$65536$ & $4127474$ & $3132$ & $3214$ & $4932$ & $2680$ & $9848$\\
$262144$ & $16508741$ & $4352$ & $827$ & $11392$ & $2560$ & $14848$\\
\bottomrule
\end{tabular}
\end{center}
At the final endpoint,
\[
 (H_{k,0})_{k=1}^8=
 (106237,84569,65768,49604,35783,24127,14429,6449),
\]
\[
 (H_{k,1})_{k=1}^8=(3,-56,-174,-146,115,-54,-111,-14),
 \qquad K_{262144}=129080.
\]
The finite Euler-product reproduction is
$0.4920202775829839485$ when both $2/\pi^2$ and the $e_k$ are replaced by
their stated products through $2^{20}$.  It is a truncated conditional
diagnostic, not an estimate establishing \eqref{eq:conditional-constant}.
All finite rows reproduce exact identities only.  No fit is used as
asymptotic evidence.

\section{Route verdict, scope, and Gates}

Route A is \texttt{GO}: the mixed hierarchy, deterministic-layer theorem,
thirteen-channel simultaneous boundary, exact two-envelope reduction, and
stationary algebraic witness are independent theorem edges.  Route B is
\texttt{STOP\_SCOPED}: the adaptive capacity is a scalar optimization of a
prefix, not an intrinsic dynamical trace, and convergence of
$(U_N+|V_N|)/N$ remains open.

The boundary is strict:
\begin{itemize}[leftmargin=*]
\item no $A_k/N$, $B_k/N$, or capacity-envelope limit is proved for the
      unresolved mixed layers.
\item the $466\to13$ computation is formal and does not establish an
      arithmetic minimal sufficient statistic.
\item the full-cancellation hypothesis and constant
      \eqref{eq:conditional-constant} are conditional and sufficient only.
\item the stationary witness is synthetic, does not have M\"obius
      squarefree statistics, and yields no M\"obius nonconvergence result.
\item no finite row, decimal Euler diagnostic, growing-modulus argument, or
      growing-clock extrapolation is promoted to an asymptotic theorem.
\end{itemize}

Gates A--E remain false/open.  There is no canonical intrinsic spectral
determinant, time-oriented scattering or unitary completion, self-adjoint
generator with an intrinsic $T\log T$ law, von Mangoldt-weighted prime-power
trace, or completed-zeta divisor equality.  This paper does not identify
Riemann zeros, construct a Hilbert--Polya operator, or prove the Riemann
Hypothesis.

\section{Conclusion}

The eight-level run hierarchy separates cleanly into deterministic
squarefree support, unconditional one-sign cancellation, and thirteen
higher mixed aggregates.  For the sixteen densities, all thirteen
aggregate limits are necessary and sufficient simultaneously.  For the
capacity, the weaker unresolved object is the single nonlinear envelope
$(U_N+|V_N|)/N$.  The stationary witness explains why low-order ternary
algebra alone cannot remove its first directional mixed layer.  A next
route therefore needs a genuine arithmetic theorem for aggregate
mixed-exponent cancellation or for the envelope itself.  Reproducing more
finite prefixes does not cross that boundary.

\section*{Reproducibility and declarations}

\textbf{Data availability.}  All source, exact outputs, tests, schema, and
SHA-256 manifests are contained in the RH-377 repository directory.

\textbf{Ethics.}  The study uses no human participants, animals, or personal
data.  No ethics approval was required.

\textbf{Author contributions.}  The RH research program performed
conceptualization, formal analysis, software, validation, and writing.

\textbf{Funding and conflicts.}  No external funding or conflict of interest
is declared.

\textbf{AI-assisted workflow.}  Drafting and executable checking used the
repository's documented multi-agent workflow.  Mathematical claims are
limited by source locks, proofs, independent audits, and reproducible
artifacts recorded here.

\bibliographystyle{plain}
\bibliography{references}
\end{document}
